\section{Conclusion}
\label{sec:conclusion}
We introduced sweep-valuation attention, which turns a directional filtration of a leaf into tokens
that carry the area and boundary length entering each slab, weighted by appearance, and attend to
the matching image region. In a controlled study with matched controls on three backbones, it gives
an ImageNet CNN a large and statistically robust gain in studio-to-field transfer (+16 points) and
keeps that gain without rotation augmentation, where strip pooling does not. The same study shows the
limits: the gain vanishes for a ViT and for the BioCLIP~2 foundation model, co-registered attention
is not its main source, and Euler-characteristic tokens, although they make the model sensitive to
topology, transfer worse than metric valuations. We release the leakage-audited protocol, the
corruption, cue-conflict and topology-destruction probes, and all code and run outputs.

\section*{Code and data availability}
All code, configurations, unit tests, the outputs of every training run (metrics, probe results,
per-image predictions and logs), the statistical analyses and the \LaTeX{} sources of this paper are
publicly available at \url{https://github.com/Devguru-codes/research_paper_leaf}. Every number
reported here is regenerated from those run outputs by the scripts in the repository, and the
properties stated in \cref{sec:method} are verified by its unit tests (\cref{app:tests}). The image
datasets are not redistributed; they are available from their original sources
(\cref{sec:data})~\citep{roopashree2020dataset,pushpa2023dimpsar}.